% Part: methods
% Chapter: proofs
% Section: inference-patterns

\documentclass[../../../include/open-logic-section]{subfiles}

\begin{document}

\olfileid{mth}{prf}{pat}

\olsection{استنباط دے ڈھنگ}

ثبوت وکھ وکھ استنباطاں توں بندے نیں۔ جدوں اسیں کوئی استنباط کردے آں، تاں عام طور اُتے «سو»، «ایس لئی»، یا «ایس کرکے» ورگا لفظ لا کے دس دیندے آں۔ استنباط اکثر اک یا دو اوہناں گلاں اُتے ٹکدا اے جیہڑیاں ثبوت وچ پہلاں ای ساڈے کول نیں۔ ایہ کوئی منّی ہوئی گل ہو سکدی اے یا کوئی نتیجہ جیہڑا اسیں پہلاں کڈھیا ہووے۔ گل صاف رکھن لئی اسیں ایہناں گلاں نوں نشان لا سکدے آں، تے استنباط کردیاں دس سکدے آں کہ اوہدے وچ کیہڑے ہور بیان ورتدے آں۔ استنباط وچ اکثر ایہ وضاحت وی ہُندی اے کہ ساڈا نواں نتیجہ پہلاں والیاں گلاں توں \emph{کیوں} نکلدا اے۔ استنباط دے کجھ عام ڈھنگ ثبوتاں وچ بہت واری ورتے جاندے نیں۔ ہِیٹھ اسیں کجھ ویکھاں گے۔ کجھ ڈھنگ، جیویں استقرا نال ثبوت، ذرا گنجل دار نیں (اوہناں دی گل اگے ہووے گی)۔

اسیں استنباط دے اک ڈھنگ دی گل پہلاں ای کیتی اے، یعنی تعریف کھولھنی یا لاؤنی۔ جدوں اسیں اک تعریف کھولھدے آں، تاں جیہڑی گل وچ تعریف ہون والی عبارت آؤندی اے، اوہنوں تعریف دین والی عبارت نال فیر آکھدے آں۔ مثلاً منّ لو اسیں ثبوت وچ پہلاں ای ثابت کر چکے آں کہ $D = E$ (ا)۔ تاں اسیں سیٹاں لئی $=$ دی تعریف لا کے ایہ نتیجہ کڈھ سکدے آں۔ «سو (ا) توں تعریف مطابق ہر !!{element} جیہڑا~$D$ دا اے، اوہ !!a{element}~$E$ دا وی اے، تے اُلٹے پاسے وی۔»

ایہ گل تھوڑا بھلیکھا پا سکدی اے کہ اسیں اکثر استنباط دا جواز اوہ کردیاں نہیں، اوہتوں پہلاں لکھدے آں۔ منّ لو اسیں ہن تک $D = E$ ثابت نہیں کیتا، پر کرنا چاہندے آں۔ جے $D = E$ ساڈا لوڑیندا نتیجہ اے، تاں ایہ مقصد وی تعریف لا کے فیر دس سکدے آں۔ $D = E$ ثابت کرن لئی سانوں ثابت کرنا پئے گا کہ ہر !!{element} جیہڑا~$D$ دا اے، اوہ !!a{element}~$E$ دا وی اے، تے اُلٹے پاسے وی۔ سو ساڈے ثبوت دی صورت ایہ ہووے گی۔ (ا) ثابت کرو کہ ہر !!{element} جیہڑا~$D$ دا اے، اوہ !!a{element}~$E$ دا وی اے۔ (ب) ہر !!{element} جیہڑا~$E$ دا اے، اوہ !!a{element}~$D$ دا وی اے۔ (ج) سو (ا) تے (ب) توں، $=$ دی تعریف مطابق، $D = E$۔ پر اسیں عام طور اُتے ایویں نہیں لکھدے۔ اسیں ایہو جیہی گل لکھ سکدے آں۔
\begin{quote}
اسیں $D = E$ وکھاؤنا چاہندے آں۔~$=$ دی تعریف مطابق ایہ وکھاؤنا اے کہ ہر !!{element} جیہڑا~$D$ دا اے، اوہ !!a{element}~$E$ دا وی اے، تے اُلٹے پاسے وی۔

(ا) \dots (ثبوت کہ ہر !!{element} جیہڑا~$D$ دا اے، اوہ !!a{element}~$E$ دا وی اے) \dots

(ب) \dots (ثبوت کہ ہر !!{element} جیہڑا~$E$ دا اے، اوہ !!a{element}~$D$ دا وی اے) \dots
\end{quote}

\subsection{عطف ورتنا}

شاید استنباط دا سبھ توں سوکھا ڈھنگ ایہ اے کہ عطف دے دو جزاں وچوں اک نوں نتیجے وجوں کڈھ لئیے۔ یعنی جے اسیں منّیا یا پہلاں ای ثابت کیتا اے کہ $p$ تے~$q$، تاں اسیں~$p$ کڈھ سکدے آں (تے~$q$ وی)۔ ایہ اینا بنیادی استنباط اے کہ اکثر ایس دا ذکر وی نہیں ہُندا۔ مثلاً $D = E$ دی تعریف کھولھ کے اسیں ثابت کر لیا کہ ہر !!{element} جیہڑا~$D$ دا اے، اوہ !!a{element}~$E$ دا وی اے، تے اُلٹے پاسے وی۔ ایس توں نتیجہ نکلدا اے کہ ہر !!{element} جیہڑا~$E$ دا اے، اوہ !!a{element}~$D$ دا وی اے (ایہو «اُلٹے پاسے وی» والا حصہ اے)۔

\subsection{عطف ثابت کرنا}

کدی جیہڑی گل تہانوں ثابت کرن لئی آکھن گے، اوہ عطف دی صورت وچ ہووے گی، «$p$ تے $q$ ثابت کرو»۔ ایس صورت وچ بس دو کم کرنے نیں۔ $p$ ثابت کرو تے فیر $q$ ثابت کرو۔ تسی ثبوت دو حصیاں وچ ونڈ سکدے او تے صفائی لئی اوہناں نوں نشان لا سکدے او۔ پہلاں نوٹ بناندیاں ورق دے اُتے «(۱) $p$ ثابت کرو» تے ورق دے وچکار «(۲) $q$ ثابت کرو» لکھ سکدے او۔ (ہو سکدا اے تہانوں عطف ثابت کرن لئی صاف نہ آکھیا ہووے پر ثبوت لئی عطف ثابت کرنا پئے۔ مثلاً جے $D = E$ ثابت کرن لئی آکھیا اے، تاں~$=$ دی تعریف کھولھ کے تہانوں ثابت کرنا پئے گا کہ ہر !!{element} جیہڑا~$D$ دا اے، اوہ !!a{element}~$E$ دا وی اے، \emph{تے} ہر !!{element} جیہڑا~$E$ دا اے، اوہ !!a{element}~$D$ دا وی اے۔)

\subsection{فصل ثابت کرنا}

جدوں جیہڑی گل تسی ثابت کردے او اوہ فصل دی صورت وچ ہووے (یعنی «$p$ یا $q$» دی صورت دا بیان ہووے)، تاں فصل دے جزاں وچوں اک نوں سچ وکھاؤنا کافی اے۔ پر ایہ لگ بھگ کدی نہیں ہُندا کہ کوئی اک جز تہاڈے قضیے دیاں منّیاں شرطاں توں سِدھا ای نکل آوے۔ اکثر قضیے دیاں منّیاں شرطاں آپ فصل دی صورت وچ ہُندیاں نیں، یا تسی وکھاؤندے او کہ اک قسم دیاں ساریاں شےواں وچ دو خاصیتاں وچوں اک اے، پر کجھ شےواں وچ پہلی تے ہورناں وچ دوجی خاصیت اے۔ ایتھے صورتاں وکھ کر کے ثبوت دین دا ڈھنگ کم آؤندا اے (ہِیٹھ ویکھو)۔

\subsection{شرطیہ ثبوت}

کئی قضیے جیہڑے تسی ویکھو گے شرطیہ صورت وچ ہون گے (یعنی وکھاؤ کہ جے $p$ پورا ہووے تاں $q$ وی سچ اے)۔ ایہناں صورتاں وچ ثبوت دی بناوٹ سوکھی اے۔ بس شرطیے دا مقدم منّ لو (ایس صورت وچ $p$) تے اوہتوں نتیجہ~$q$ ثابت کرو۔ سو جے قضیہ ایہ آکھدا اے، «جے $p$ تاں $q$»، تاں ثبوت «$p$ منّ لو» توں شروع کرو تے آخر تک~$q$ ثابت کر لو۔

شرطیہ کئی طرحاں نال دَسے جا سکدے نیں۔ سو «جے $p$ تاں $q$» دی تھاں قضیہ آکھ سکدا اے «$p$ صرف اُدوں جے $q$»، «$q$ جے $p$»، یا «$q$، بشرطیکہ $p$»۔ ایہناں ساریاں دا مطلب اکّو اے، تے ہر واری $p$ منّ کے اوس منّی ہوئی گل توں~$q$ ثابت کرنا پئے گا۔ چیتے کرو کہ دوپاسا شرطیہ («$p$ اُدوں تے صرف اُدوں جدوں $q$») اصل وچ دو شرطیاں دا جوڑ اے، جے $p$ تاں $q$، تے جے $q$ تاں~$p$۔ سو تہانوں بس دو شرطیہ ثبوت دینے پَین گے، اک پہلے شرطیے لئی تے اک دوجے لئی۔ پر کدی «اُدوں تے صرف اُدوں جدوں» والا بیان ایویں وی ثابت ہو سکدا اے کہ ایس طرح دے کئی ہور بیان اک لڑی وچ جوڑ دیو، «$p$» توں شروع کرو تے «$q$» اُتے مُکو۔ پر اوس صورت وچ پکا کر لو کہ ہر قدم واقعی «اُدوں تے صرف اُدوں جدوں» والا اے۔

\subsection{ساریاں لئی دعوے}

ساریاں لئی کوئی دعویٰ ورتنا سوکھا اے۔ جے کوئی گل ہر شے لئی سچ اے، تاں ہر خاص شے لئی وی سچ اے۔ سو مثلاً جے تہاڈے ثبوت دی مفروضہ شرط $A \subseteq B$ اے، تاں (تعریف~$\subseteq$ کھولھ کے) ایہدا مطلب اے کہ ہر $x \in A$ لئی، $x \in B$۔ سو جے تہانوں پہلاں پتا اے کہ $z \in A$، تاں نتیجہ کڈھ سکدے او کہ $z \in B$۔

ساریاں لئی کوئی دعویٰ ثابت کرنا ذرا اوکھا لگ سکدا اے۔ عام طور اُتے ایہ بیان ایہو جیہی صورت وچ ہُندے نیں۔ «جے $x$ وچ~$P$ اے، تاں اوہدے وچ~$Q$ وی اے» یا «سارے $P$، $Q$ نیں»۔ بے شک ہو سکدا اے ایہ عین ایس صورت وچ نہ ہووے، تے عین کیہڑی گل ثابت کرن لئی آکھی گئی اے ایہ سمجھنا تھوڑی مشق منگدا اے۔ پر سانوں اکثر ثابت کرنا پیندا اے کہ اک خاصیت والیاں ساریاں شےواں وچ اک ہور خاصیت وی اے۔

ساریاں لئی دعویٰ ثابت کرن دا طریقہ ایہ اے کہ اوہناں شےواں لئی ناں یا متغیر پیش کرو جیہناں وچ پہلی خاصیت اے، تے فیر وکھاؤ کہ اوہناں وچ دوجی خاصیت وی اے۔ ایہ گل ایویں آکھ سکدے آں کہ \emph{سارے}~$P$ لئی کوئی گل ثابت کرن واسطے اک \emph{من مانے}~$P$ لئی اوہ گل ثابت کرنی پئے گی۔ جیہڑا ناں پیش کیتا اے، اوہ اک من مانے~$P$ دا ناں اے۔ من مانیاں شےواں دے ناں لئی اسیں عام طور اُتے اکّو اکّو اکھر ورتدے آں، تے اکھراں دی چُون اکثر رواج مطابق ہُندی اے۔ مثلاً قدرتی عدداں لئی $n$، !!{formula}s لئی $!A$، سیٹاں لئی $A$، فنکشناں لئی $f$، وغیرہ۔

کم ایہ اے کہ سارے ثبوت وچ عمومیت قائم رکھو۔ شروع وچ منّ لو کہ اک من مانی شے («$x$») وچ خاصیت~$P$ اے، تے وکھاؤ (بس تعریفاں یا اوہناں گلاں توں جیہڑیاں تہانوں منّن دی کھُل اے) کہ $x$ وچ خاصیت~$Q$ اے۔ کیوں جے تسی $x$ بارے خاص طور اُتے کجھ نہیں دسیا، صرف ایہ کہ اوہدے وچ خاصیت $P$ اے، تسی آکھ سکدے او کہ ہر $P$ والی شے وچ خاصیت~$Q$ اے۔ مختصر گل ایہ اے کہ $x$ خاصیت~$P$ والیاں \emph{ساریاں} شےواں دی نمائندگی کردا اے۔

\begin{prop}
  سارے سیٹاں $A$ تے $B$ لئی، $A \subseteq A \cup B$۔
\end{prop}

\begin{proof}
  $A$ تے $B$ من مانے سیٹ ہون۔ اسیں وکھاؤنا چاہندے آں کہ $A
  \subseteq A \cup B$۔ $\subseteq$ دی تعریف مطابق ایہ وکھاؤنا اے کہ ہر $x$ لئی، جے $x \in A$ تاں $x \in A \cup B$۔ سو $x \in A$، سیٹ~$A$ دا اک من مانا !!{element} ہووے۔ سانوں وکھاؤنا اے کہ $x \in A
  \cup B$۔ کیوں جے $x \in A$، ایس لئی $x \in A$ یا $x \in B$۔ سو $x \in
  \Setabs{x}{x \in A \lor x \in B}$۔ پر $\cup
  $ دی تعریف مطابق ایہدا مطلب $x \in A \cup B$ اے۔
\end{proof}

\subsection{صورتاں وکھ کر کے ثبوت}

منّ لو فصل تہاڈی منّی ہوئی گل اے یا پہلاں ای ثابت کیتا نتیجہ اے، یعنی تسی منّیا یا ثابت کیتا اے کہ $p$ یا $q$ سچ اے۔ تسی $r$ ثابت کرنا چاہندے او۔ ایہ دو قدماں وچ کرو۔ پہلاں $p$ نوں سچ منّ کے~$r$ ثابت کرو۔ فیر $q$ نوں سچ منّ کے~$r$ فیر ثابت کرو۔ ایہ طریقہ ایس لئی چلدا اے کہ اسیں منّیا اے یا سانوں پتا اے کہ دوناں وچوں اک صورت پوری ہُندی اے۔ دو قدم ثابت کردے نیں کہ دوہاں وچوں کوئی وی صورت~$r$ دے سچ ہون لئی کافی اے۔ (جے دوویں سچ نیں، تاں $r$~دے سچ ہون دی اک دی بجائے دو وجوہاں نیں۔ $r$~دا سچ ہونا، $p$ تے~$q$ دوویں منّ کے وکھرا ثابت کرن دی لوڑ نہیں۔) اپنا طریقہ دسّن لئی اسیں آکھدے آں کہ «اسیں صورتاں وکھ کردے آں»۔ مثلاً منّ لو سانوں پتا اے کہ $x \in B \cup C$۔ $B \cup C$ دی تعریف $\Setabs{x}{x \in B \text{ یا } x \in C}$ اے۔ دوجے لفظاں وچ تعریف مطابق $x \in B$ یا $x \in C$۔ ایس توں $x \in A$ ثابت کرن لئی پہلاں $x \in B$ منّو، تے ایس منّی ہوئی گل توں $x \in A$ ثابت کرو۔ فیر $x \in C$ منّو، تے ایس توں فیر $x \in A$ ثابت کرو۔ منّی ہوئی گل تھلّے «اسیں صورتاں وکھ کردے آں» لکھو، اوہدے تھلّے «صورت (۱) $x \in B$»، تے ورق دے وچکار «صورت (۲) $x \in C$» لکھو۔ فیر ورق دا اُتلا تے تھلّلا ادھ بھرو۔

صورتاں وکھ کر کے ثبوت اوہدوں خاص کم آؤندا اے جدوں ثابت کرن والی گل آپ فصل دی صورت وچ ہووے۔ اک سوکھی مثال ایہ اے۔

\begin{prop}
منّ لو $B \subseteq D$ تے $C \subseteq E$۔ تاں $B \cup C \subseteq
D \cup E$۔
\end{prop}

\begin{proof}
  منّ لو (ا) $B \subseteq D$ تے (ب) $C \subseteq E$۔ تعریف مطابق کوئی وی $x \in B$، $\in D$ وی اے (ج)، تے کوئی وی $x \in C$، $\in E$ وی اے (د)۔ $B \cup C \subseteq D \cup E$ وکھاؤن لئی سانوں وکھاؤنا پئے گا کہ جے $x \in B \cup C$ تاں $x \in D \cup E$ ($\subseteq$ دی تعریف مطابق)۔ $x \in B \cup C$ اُدوں تے صرف اُدوں جدوں $x \in B$ یا $x \in
  C$ (تعریف~$\cup$ مطابق)۔ ایویں ای $x \in D \cup E$ اُدوں تے صرف اُدوں جدوں $x \in
  D$ یا $x \in E$۔ سو سانوں ایہ وکھاؤنا اے کہ کسے وی $x$ لئی، جے $x \in B$ یا $x \in C$، تاں $x \in D$ یا $x \in E$۔

  \begin{quote}
  ہن تک اسیں صرف تعریفاں کھولھیاں نیں!{} اسیں اپنا قضیہ $\subseteq$ تے $\cup$ توں بناں فیر بیان کیتا اے، تے ہن ساریاں لئی اک شرطیہ دعویٰ ثابت کرنا رہ گیا اے۔ اُتلی گل مطابق، ایہ کم کرن لئی $x$ نوں اک ایہو جیہی شے منّو جیہدے بارے اسیں «جے» والا حصہ سچ منّدے آں۔ اسیں اگے وکھاواں گے کہ «تاں» والا حصہ وی سچ اے۔ دوجے لفظاں وچ، اسیں منّاں گے کہ $x
  \in B$ یا $x \in C$، تے وکھاواں گے کہ $x \in D$ یا $x \in
  E$۔\footnote{ایہ پیرا بس سمجھاؤندا اے کہ اسیں کی کردے آں۔ ایہ ثبوت دا حصہ نہیں، تے اپنے ثبوت لکھدیاں تہانوں اینی تفصیل وچ جان دی لوڑ نہیں۔}
  \end{quote}

  منّ لو $x \in B$ یا $x \in C$۔ سانوں وکھاؤنا اے کہ $x \in D$ یا $x \in E$۔ اسیں صورتاں وکھ کردے آں۔

  صورت ۱۔ $x \in B$۔ (ج) مطابق $x \in D$۔ سو $x \in D$ یا $x \in
  E$۔ (ایتھے اسیں اوہ استنباط کیتا اے جیہدی گل پچھلے ذیلی حصے وچ کیتی سی!)

  صورت ۲۔ $x \in C$۔ (د) مطابق $x \in E$۔ سو $x \in D$ یا $x \in E$۔
 \end{proof}

\subsection{موجودگی دا دعویٰ ثابت کرنا}

جدوں موجودگی دا دعویٰ ثابت کرن لئی آکھیا جائے، تاں سوال عام طور اُتے ایس صورت وچ ہووے گا۔ «ثابت کرو کہ اک~$x$ اے جیہدے لئی $\dots x \dots$»، یعنی کوئی شے اے جیہدے وچ «$\dots x
\dots$» نال دسی خاصیت اے۔ ایس صورت وچ تہانوں کوئی ڈھکواں شے لبھنی تے وکھاؤنا پئے گا کہ اوہدے وچ لوڑیندی خاصیت اے۔ گل سِدھی لگدی اے، پر ایس قسم دا ثبوت اوکھا ہو سکدا اے۔ عام طور اُتے اک شے \emph{بنانی} یا \emph{تعریف کرنی} تے ثابت کرنا پیندا اے کہ ایویں تعریف کیتی شے وچ لوڑیندی خاصیت اے۔ ٹھیک شے لبھنی اوکھی ہو سکدی اے، اوہدے وچ لوڑیندی خاصیت ثابت کرنی اوکھی ہو سکدی اے، تے کدی ایہ وکھاؤنا وی اوکھا ہو سکدا اے کہ تسی واقعی کوئی شے تعریف کر لئی اے!{}

عام طور اُتے تسی شے دس کے ایہ کم لکھو گے، مثلاً «$x$، \dots ہووے» (جتھے \dots{} دسدا اے کہ تہاڈے ذہن وچ کیہڑی شے اے)۔ ہو سکدا اے پہلاں ثابت کرنا پئے کہ $\dots$ واقعی اک موجود شے دا بیان اے، تے فیر وکھاؤ کہ $x$ وچ خاصیت~$Q$ اے۔ اک سوکھی مثال ایہ اے۔

\begin{prop}
  منّ لو $x \in B$۔ تاں اک~$A$ اے جیہدے لئی $A \subseteq
  B$ تے $A \neq \emptyset$۔
\end{prop}

\begin{proof}
  $x \in B$ منّ لو۔ $A = \{x\}$ ہووے۔
  \begin{quote}
    ایتھے اسیں سیٹ~$A$ دے !!{element}s گِنوا کے اوہدی تعریف دتی اے۔ کیوں جے اسیں منّیا اے کہ $x$ اک شے اے، تے اسیں !!{element}s گِنوا کے ہمیشہ اک سیٹ بنا سکدے آں، سانوں ایہ وکھاؤن دی لوڑ نہیں کہ ایتھے اسیں واقعی اک سیٹ~$A$ تعریف کیتا اے۔ پر سانوں ہن وی وکھاؤنا اے کہ $A$ وچ قضیے دیاں لوڑیندیاں خاصیتاں نیں۔ ایس توں بناں ثبوت پورا نہیں!{}
  \end{quote}
  کیوں جے $x \in A$، ایس لئی $A \neq \emptyset$۔
  \begin{quote}
    ایہ $A$ دی $\{x\}$ وجوں تعریف تے ایہناں صاف گلاں اُتے ٹکدا اے کہ $x \in \{x\}$ تے $x \notin \emptyset$۔
  \end{quote}
  کیوں جے $x$ اکّو اک !!{element} اے~$\{x\}$ دا، تے $x \in B$، سو ہر !!{element} جیہڑا~$A$ دا اے، اوہ !!a{element}~$B$ دا وی اے۔~$\subseteq$ دی تعریف مطابق $A \subseteq B$۔
\end{proof}

\subsection{موجودگی دے دعوے ورتنے}

منّ لو تہانوں پتا اے کہ موجودگی دا کوئی دعویٰ سچ اے (تسی اوہنوں ثابت کر چکے او یا ایہ اک منّی ہوئی شرط اے جیہڑی تسی ورت سکدے او)، مثلاً «کسے~$x$ لئی، $x \in A$» یا «اک $x \in A$ اے»۔ جے تسی اینوں ثبوت وچ ورتنا چاہندے او، تاں منّ سکدے او کہ تہاڈے کول اوہناں شےواں وچوں اک دا ناں اے جیہڑیاں تہاڈی شرط مطابق موجود نیں۔ کیوں جے $A$ وچ گھٹ توں گھٹ اک شے اے، ایہو جیہیاں شےواں نیں جیہناں ول اوہ ناں اشارہ کر سکدا اے۔ ہو سکدا اے تسی اوہناں وچوں کوئی اک شے چُن نہ سکو یا اوہدا ہور بیان نہ دے سکو (بس اینا کہ اوہ $\in A$ اے)۔ پر ثبوت لئی منّ سکدے او کہ تسی اوہ شے چُن لئی اے تے اوہدا ناں رکھ سکدے او۔ ضروری اے کہ اوہ ناں چُنو جیہڑا تسی پہلاں نہیں ورتیا تے جیہڑا تہاڈیاں منّیاں شرطاں وچ نہیں آؤندا، نہیں تاں گل وگڑ سکدی اے۔ ثبوت وچ «کسے $x$ لئی، $x \in A$» توں «$a
\in A$ ہووے» ول جا کے ایہ دس دیو۔ ہُن تسی~$a$ بارے دلیل دے سکدے او، ہور منّیاں شرطاں ورت سکدے او، وغیرہ، جدوں تک تسی نتیجے $p$ تک نہ پہنچ جاؤ۔ جے $p$ وچ ہن~$a$ دا ذکر نہیں، تاں $p$ ایس منّی ہوئی گل توں آزاد اے کہ $a \in A$، تے تسی وکھا دتا اے کہ اوہ بس ایس منّی ہوئی گل توں نکلدا اے کہ «کسے $x$ لئی، $x \in A$»۔

\begin{prop}
جے $A \neq \emptyset$، تاں $A \cup B \neq \emptyset$۔
\end{prop}

\begin{proof}
  منّ لو $A \neq \emptyset$۔ سو کسے $x$ لئی، $x \in A$۔
  \begin{quote}
    ایتھے اسیں پہلاں بس قضیے دی مفروضہ شرط فیر دسی اے۔ ایس شرط، یعنی $A \neq \emptyset$، وچ موجودگی دا اک دعویٰ لُکیا اے، جیہڑا کجھ تعریفاں کھولھ کے ای لبھدا اے۔ $=$ دی تعریف دسدی اے کہ $A = \emptyset$ اُدوں تے صرف اُدوں جدوں ہر $x \in A$، $\in \emptyset$ وی اے، تے ہر $x \in
    \emptyset$، $\in A$ وی اے۔ دوہاں پاسیاں دی نفی لئیے تاں $A \neq
    \emptyset$ اُدوں تے صرف اُدوں جدوں یا کوئی $x \in A$، $\notin \emptyset$ اے، یا کوئی $x \in \emptyset$، $\notin A$ اے۔ کیوں جے کوئی شے $\in \emptyset$ نہیں، فصل دا دوجا جز کدی سچ نہیں ہو سکدا، تے «$x \in A$ تے $x
    \notin \emptyset$» بس $x \in A$ رہ جاندا اے۔ سو $x \neq
    \emptyset$ اُدوں تے صرف اُدوں جدوں کسے $x$ لئی، $x \in A$۔ ایہ موجودگی دا دعویٰ اے۔ ہُن اسیں~$A$ دے !!{element}s وچوں اک دا ناں پیش کر کے اوہ دعویٰ ورتدے آں:
  \end{quote}
  $a \in A$ ہووے۔
  \begin{quote}
    ہُن اسیں اوہناں شےواں وچوں اک دا ناں پیش کیتا اے جیہڑیاں~$\in A$ نیں۔ اسیں اگے~$a$ بارے دلیل دیواں گے، پر دھیان رکھاں گے کہ بس $a \in A$ منّیئے، تے ہور کجھ نہ منّیئے:
  \end{quote}
  کیوں جے $a \in A$، ایس لئی $a \in A \cup B$، تعریف~$\cup$ مطابق۔ سو کسے $x$ لئی، $x \in A \cup B$، یعنی $A \cup B \neq \emptyset$۔
  \begin{quote}
    ایس آخری قدم وچ اسیں «$a \in A \cup B$» توں «کسے $x$ لئی، $x \in A \cup B$» تک گئے آں۔ ایس وچ ہن $a$ دا ذکر نہیں، سو سانوں پتا اے کہ «کسے $x$ لئی، $x \in A \cup B$» بس «کسے $x$ لئی، $x \in A$» توں نکلدا اے۔ پر ایہدا مطلب اے کہ $A \cup B \neq
    \emptyset$۔
  \end{quote}
\end{proof}

شاید چنگا طریقہ اے کہ «$x$» ورگے بندھے متغیر، $a$ ورگے منّے ہوئے ناںواں توں وکھ رکھو، جیویں اسیں کیتا۔ پر کم کردیاں اسیں اکثر ایہ نہیں کردے تے بس $x$ ای ورتدے آں، ایویں:
\begin{quote}
منّ لو $A \neq \emptyset$، یعنی اک $x \in A$ اے۔ $\cup$ دی تعریف مطابق $x \in A \cup B$۔ سو $A \cup B \neq \emptyset$۔
\end{quote}
پر جدوں تسی ایویں کردے او، تاں اچھی طرح دھیان رکھنا پئے گا کہ وکھ وکھ موجودگی دے دعویاں لئی وکھرے $x$ تے $y$ ورتو۔ مثلاً ہِیٹھلا «جے $A \neq
\emptyset$ تے $B \neq \emptyset$ تاں $A \cap B \neq \emptyset$» دا ٹھیک ثبوت \emph{نہیں} (تے ایہ دعویٰ سچ نہیں)۔
\begin{quote}
منّ لو $A \neq \emptyset$ تے $B \neq \emptyset$۔ سو کسے $x$ لئی، $x
\in A$، تے کسے $x$ لئی، $x \in B$ وی۔ کیوں جے $x \in A$ تے $x \in
B$، ایس لئی $x \in A \cap B$، تعریف~$\cap$ مطابق۔ سو $A \cap B \neq
\emptyset$۔
\end{quote}
کی تسی غلط قدم لبھ کے سمجھا سکدے او کہ نتیجہ کیوں نہیں نکلدا؟

\end{document}
